% ============================================================
\subsection{Empirical recurrent-scenario mining}
\label{sec:chap4_scenario_mining}
% ============================================================

The accident-level analysis first aims to identify configurations of semantic
factors that recur across accident narratives and that can therefore support
the investigation of recurrent accident processes. Scenario discovery is
performed directly from the accident--factor matrix and independently of the
Bayesian-network model introduced subsequently. This separation is deliberate:
recurrence is treated as an empirical property of the observed accident
narratives, whereas the Bayesian network addresses the conditional organization
of the factors and the extent to which its joint distribution reproduces the
frequency of the recurrent configurations.

The analysis is formulated as a constrained frequent-itemset mining problem.
In classical frequent-pattern mining, each observation is represented as a
transaction containing a set of items, and recurrent combinations are
identified according to their support in the transaction database
\citep{agrawal1993mining}. Here, the transactions are accident narratives and
the items are the retained semantic factors. The candidate space is restricted
a priori according to the predefined accident-process roles, following the
general principle of constraint-based pattern mining
\citep{ng1998exploratory}. These restrictions focus the search on parsimonious
configurations that retain an interpretable accident-process structure for
subsequent prevention-oriented analysis.


% ------------------------------------------------------------
\subsubsection{Admissible scenarios and empirical recurrence}
\label{sec:chap4_scenario_candidates}
% ------------------------------------------------------------

For accident narrative \(i\), let
\[
\mathcal T_i
=
\left\{
j\in\{1,\ldots,p\}:X_{ij}=1
\right\}
\]
denote the transaction containing all retained semantic factors identified in
that narrative. A candidate scenario contains one or two upstream factors from
roles \roleAzero{} and/or \roleAone{}, exactly one event/deviation factor from
role \roleB{}, and exactly one consequence factor from role \roleC{}. The
admissible scenario family is therefore

\[
\mathfrak S
=
\left\{
\mathcal F^{\mathrm{up}}\cup\{b,c\} :
\mathcal F^{\mathrm{up}}
\subseteq
\mathcal J_{\mathrm{A0}}\cup\mathcal J_{\mathrm{A1}},
\;
1\leq|\mathcal F^{\mathrm{up}}|\leq2,
\;
b\in\mathcal J_{\mathrm B},
\;
c\in\mathcal J_{\mathrm C}
\right\}.
\]

Restricting the upstream component to at most two factors yields configurations
centered on a specific accident event and consequence while still allowing
combined contextual or adverse-condition factors to be represented. When two
upstream factors are included, they may belong either to the same role or to
different roles.

For scenario \(s\), let
\(\mathcal S_s=\mathcal F_s^{\mathrm{up}}\cup\{b_s,c_s\}\).
Narrative \(i\) contains the scenario whenever
\(\mathcal S_s\subseteq\mathcal T_i\). Factors outside
\(\mathcal S_s\) remain unrestricted, so scenario occurrence corresponds to
itemset containment rather than equality with the complete accident profile.
The recurrence count and empirical support are

\[
n_s
=
\sum_{i=1}^{N_A}
\mathbb I
\left\{
\mathcal S_s\subseteq\mathcal T_i
\right\},
\qquad
\widehat{\operatorname{supp}}(\mathcal S_s)
=
\frac{n_s}{N_A}.
\]

A configuration is considered empirically recurrent when
\(n_s\geq n_{\min}\). The threshold \(n_{\min}\) is a recurrence criterion and
not a statistical-significance threshold. Its sensitivity is examined in the
empirical application to determine whether the principal recurrent
configurations depend strongly on the selected minimum count.


% ------------------------------------------------------------
\subsubsection{Association with consequences and redundancy reduction}
\label{sec:chap4_scenario_measures}
% ------------------------------------------------------------

Recurrence quantifies how often the complete configuration is observed, but it
does not characterize whether its consequence is specifically enriched in the
corresponding upstream accident configuration. For scenario \(s\), define the
association-rule antecedent as
\(\mathcal A_s=\mathcal F_s^{\mathrm{up}}\cup\{b_s\}\), so that the rule of
interest is \(\mathcal A_s\Rightarrow c_s\). Let \(n_s^A\) denote the number
of narratives containing \(\mathcal A_s\), and let

\[
\widehat p_{c_s}
=
\frac{1}{N_A}
\sum_{i=1}^{N_A}X_{i,c_s}
\]

denote the marginal prevalence of consequence \(c_s\). Following standard
association-rule terminology
\citep{agrawal1993mining,brin1997dynamic}, confidence and lift are defined as

\[
\widehat{\operatorname{conf}}_s
=
\frac{n_s}{n_s^A},
\qquad
\widehat{\operatorname{lift}}_s
=
\frac{\widehat{\operatorname{conf}}_s}
     {\widehat p_{c_s}},
\]

provided that \(n_s^A>0\). Confidence estimates the frequency of consequence
\(c_s\) among narratives containing the antecedent, whereas lift compares this
conditional frequency with the marginal prevalence of the same consequence.
A lift greater than one therefore indicates enrichment of \(c_s\) within the
considered accident configuration.

Recurrence and association are kept conceptually separate. A scenario may be
recurrent without exhibiting strong consequence enrichment, while a
configuration with high lift may occur in only a small number of accidents.
Confidence and lift are consequently used to characterize retained scenarios
but do not determine whether a configuration satisfies the recurrence
criterion.

Redundancy among recurrent configurations is reduced by adapting the
closed-itemset principle \citep{pasquier1999discovering} to the predefined
admissible family \(\mathfrak S\). A recurrent scenario
\(\mathcal S_s\) is retained as closed within this family when

\[
\nexists\,
\mathcal S_{s'}\in\mathfrak S
\quad\text{such that}\quad
\mathcal S_s\subsetneq\mathcal S_{s'}
\quad\text{and}\quad
n_{s'}=n_s.
\]

Thus, a smaller configuration is removed when an admissible strict superset
occurs in exactly the same accident narratives. Closure is used here to
reduce redundancy in the description of recurrence. Since it is defined
within the restricted scenario family, it is not interpreted as a
lossless representation of all possible association measures outside
\(\mathfrak S\).


% ============================================================
\subsection{Role-constrained Bayesian modeling of accident processes}
\label{sec:chap4_global_bn}
% ============================================================

Recurrent-scenario mining identifies semantic factors that repeatedly occur
together, but joint occurrence alone does not determine whether the
relationships among these factors remain informative once other factors are
taken into account. A Bayesian network is therefore learned over the same
accident-level binary variables to characterize their conditional dependence
structure.

The network serves two complementary purposes. First, it identifies direct
conditional dependencies that are retained after accounting for other selected
parents under the predefined accident-process organization. This provides a
more selective characterization of the relationships occurring within
recurrent scenarios than co-occurrence alone. Second, because the fitted
network defines a joint probability distribution over all retained factors, it
provides a probabilistic background model against which the empirical support
of the recurrent scenarios can subsequently be compared.

The network is not interpreted as a causal model. Edge directions are
restricted a priori according to the accident-process roles and therefore
represent conditional dependencies within this predefined structural
framework rather than causal effects inferred from the observational data.


% ------------------------------------------------------------
\subsubsection{Bayesian-network factorization and accident-process constraints}
\label{sec:chap4_global_bn_factorization}
% ------------------------------------------------------------

Let
\(\mathbf X=(X_1,\ldots,X_p)\)
denote the \(p\) binary variables corresponding to the retained semantic
factors. For each factor \(X_j\), let
\(\rho(j)\in\mathcal R\)
denote its accident-process role, and let

\[
\mathcal J_r
=
\{j:\rho(j)=r\}
\]

denote the set of factor indices associated with role \(r\).

A Bayesian network consists of a directed acyclic graph
\(G=(V,E)\), with \(V=\{X_1,\ldots,X_p\}\), together with a collection of
local conditional probability distributions. If
\(\operatorname{Pa}_G(j)\) denotes the set of indices of the direct parents
of \(X_j\), the joint distribution factorizes according to the graph as

\[
P(\mathbf X=\mathbf x\mid G,\Theta_G)
=
\prod_{j=1}^{p}
P\left(
X_j=x_j
\mid
\mathbf X_{\operatorname{Pa}_G(j)}
=
\mathbf x_{\operatorname{Pa}_G(j)}
\right).
\]

Since all variables are binary, the local distribution of \(X_j\) for a
parent configuration \(\pi\) is characterized by the Bernoulli parameter

\[
\theta_{j,\pi}
=
P\left(
X_j=1
\mid
\mathbf X_{\operatorname{Pa}_G(j)}=\pi
\right).
\]

The Bayesian-network structure is not learned over the unrestricted space of
all DAGs. Structural constraints can be incorporated into score-based
Bayesian-network learning to restrict the set of candidate graphs according
to prior or expert knowledge
\citep{decampos2009structure,li2018constraints,
ramaswamy2022expertconstraints}. In the present analysis, these restrictions
are derived from the predefined organization of the accident process. The
admissible ordered role pairs are

\[
\mathcal A_{\mathrm{role}}
=
\left\{
(\mathrm{A0},\mathrm{A1}),
(\mathrm{A0},\mathrm{B}),
(\mathrm{A1},\mathrm{B}),
(\mathrm{B},\mathrm{C})
\right\}.
\]

Accordingly, A0 factors have no admissible parents; A1 factors may have
parents only among A0 factors; B factors may have parents among A0 and A1
factors; and C factors may have parents only among B factors. Intra-role
edges, reverse-process edges, direct A0-to-C and A1-to-C edges, and edges
originating from consequence factors are excluded.

For each node \(X_j\), let \(\mathcal U_j\) denote its set of admissible
potential parents:

\[
\mathcal U_j
=
\begin{cases}
\varnothing,
& \rho(j)=\mathrm{A0},
\\[0.3em]
\mathcal J_{\mathrm{A0}},
& \rho(j)=\mathrm{A1},
\\[0.3em]
\mathcal J_{\mathrm{A0}}\cup\mathcal J_{\mathrm{A1}},
& \rho(j)=\mathrm{B},
\\[0.3em]
\mathcal J_{\mathrm{B}},
& \rho(j)=\mathrm{C}.
\end{cases}
\]

The admissible parent sets for node \(X_j\) are

\[
\mathcal P_j
=
\left\{
P\subseteq\mathcal U_j:
|P|\leq d_{\max}
\right\},
\qquad
d_{\max}=2.
\]

If \(m_j=|\mathcal U_j|\), the number of candidate parent sets considered for
node \(X_j\) is therefore

\[
|\mathcal P_j|
=
\sum_{d=0}^{\min(d_{\max},m_j)}
\binom{m_j}{d}.
\]

The indegree bound has both statistical and computational motivations. For a
binary child variable with \(d\) binary parents, the conditional-probability
table contains \(2^d\) parent configurations and therefore \(2^d\) free
Bernoulli parameters. Restricting the number of parents reduces the risk of
very sparsely populated conditional cells while still allowing a downstream
factor to depend directly on more than one upstream factor.

The role constraints also induce the strict ordering

\[
\mathrm{A0}
\prec
\mathrm{A1}
\prec
\mathrm{B}
\prec
\mathrm{C}.
\]

Every admissible edge progresses according to this ordering. A directed path
can therefore only move from an earlier role towards a later role and cannot
return to its starting node. Acyclicity is consequently guaranteed by the
definition of the admissible graph space.

\begin{figure}[!htbp]
\centering

\begin{tikzpicture}[
    node distance=0.9cm,
    >=Latex,
    base/.style={
        rectangle,
        rounded corners=2mm,
        minimum width=2.6cm,
        minimum height=1.05cm,
        align=center,
        font=\sffamily\small,
        line width=1pt,
        text=black
    },
    nodeA0/.style={
        base,
        fill=A0fill,
        draw=A0border
    },
    nodeA1/.style={
        base,
        fill=A1fill,
        draw=A1border
    },
    nodeB/.style={
        base,
        fill=Bfill,
        draw=Bborder
    },
    nodeC/.style={
        base,
        fill=Cfill,
        draw=Cborder
    },
    arrow/.style={
        ->,
        line width=1.05pt,
        draw=black!70,
        shorten >=1.5pt,
        shorten <=1.5pt
    }
]

\node[nodeA0] (A0) {
    \textbf{A0}\\[0.1em]
    {\scriptsize Work context}
};

\node[nodeA1, right=of A0] (A1) {
    \textbf{A1}\\[0.1em]
    {\scriptsize Adverse condition}
};

\node[nodeB, right=of A1] (B) {
    \textbf{B}\\[0.1em]
    {\scriptsize Event / deviation}
};

\node[nodeC, right=of B] (C) {
    \textbf{C}\\[0.1em]
    {\scriptsize Consequence}
};

\draw[arrow] (A0) -- (A1);
\draw[arrow] (A1) -- (B);
\draw[arrow] (B) -- (C);
\draw[arrow]
    (A0.south)
    to[bend right=22]
    (B.south);

\end{tikzpicture}

\caption{
Role-level architecture of the constrained Bayesian-network search space.
Arrows indicate admissible directions of direct conditional dependence
between accident-process roles and do not represent causal effects.
}
\label{fig:chap4_global_bn_architecture}
\end{figure}

The absence of an admissible edge does not imply marginal independence between
the corresponding variables. Factors within the same role may, for example,
co-occur frequently even though intra-role dependencies are not represented
as edges. More generally, the role constraints determine which direct
conditional dependencies the model is allowed to represent. They encode the
predefined temporal and conceptual organization of the accident process but
do not identify causal relationships from the observational data.


% ------------------------------------------------------------
\subsubsection{Local likelihood and decomposable BIC score}
\label{sec:chap4_global_bn_score}
% ------------------------------------------------------------

Let
\(\mathbf x_i=(x_{i1},\ldots,x_{ip})\),
\(i=1,\ldots,N_A\),
denote the binary factor profile associated with accident narrative \(i\).
Accident narratives are treated as independent observational units at this
stage of the analysis. For a candidate parent set
\(P\in\mathcal P_j\), let

\[
N_{j,\pi,x}(P)
=
\sum_{i=1}^{N_A}
\mathbb I
\left\{
\mathbf x_{i,P}=\pi,\;
x_{ij}=x
\right\},
\qquad
x\in\{0,1\},
\]

and define
\(N_{j,\pi}(P)
=
N_{j,\pi,0}(P)+N_{j,\pi,1}(P)\).
For every observed parent configuration satisfying
\(N_{j,\pi}(P)>0\), the maximum-likelihood estimate of the Bernoulli
parameter is

\[
\widehat\theta_{j,\pi}(P)
=
\frac{N_{j,\pi,1}(P)}
     {N_{j,\pi}(P)}.
\]

Using the convention \(0\log 0=0\), the corresponding maximized local
log-likelihood is

\[
\ell_j(P)
=
\sum_{\substack{
\pi\in\{0,1\}^{|P|}\\
N_{j,\pi}(P)>0}}
\left[
N_{j,\pi,1}(P)
\log\widehat\theta_{j,\pi}(P)
+
N_{j,\pi,0}(P)
\log
\left(
1-\widehat\theta_{j,\pi}(P)
\right)
\right].
\]

Parent configurations that are absent from the sample do not contribute to
the maximized likelihood. No arbitrary conditional probability is therefore
required for BIC computation.

Candidate structures are compared using the Bayesian information criterion
\citep{schwarz1978estimating}. For a binary child with \(|P|\) binary
parents, the nominal number of local free parameters is \(2^{|P|}\).
The local BIC contribution is consequently

\[
\operatorname{BIC}_j(P)
=
-2\ell_j(P)
+
2^{|P|}
\log N_A.
\]

Both the maximized log-likelihood and the complexity penalty decompose over
the nodes. Hence, for graph \(G\),

\[
\operatorname{BIC}(G)
=
\sum_{j=1}^{p}
\operatorname{BIC}_j
\left(
\operatorname{Pa}_G(j)
\right).
\]

The decomposability of score-based Bayesian-network criteria is a standard
property exploited in exact structure-learning algorithms
\citep{decampos2009structure}. It does not, by itself, imply that parent sets
can generally be optimized independently. In an unrestricted search, the
individually optimal parent sets may jointly create a directed cycle, so
acyclicity couples the local optimization problems. Exact procedures must
therefore ordinarily account for this global graph constraint. De Campos
et al.~\citep{decampos2009structure} show, however, that this additional
search disappears when a fixed ordering guarantees that every candidate
edge is acyclic. The accident-process constraints above yield precisely this
property.


% ------------------------------------------------------------
\subsubsection{Exact structure learning under the role constraints}
\label{sec:chap4_global_bn_exact}
% ------------------------------------------------------------

Let \(\mathcal G_{\mathrm{role}}\) denote the class of directed graphs whose
parent sets satisfy
\(\operatorname{Pa}_G(j)\in\mathcal P_j\)
for every \(j\). Because each admissible parent belongs to an earlier role,
all graphs in this class satisfy the required accident-process ordering.

\begin{proposition} Under the predefined role ordering and the decomposable BIC criterion, any graph obtained by independently selecting a BIC-minimizing admissible parent set for each node is globally BIC-optimal within \(\mathcal G_{\mathrm{role}}\).
\end{proposition}

\begin{proof} For each node \(X_j\), let \(\mathcal P_j\) denote its set of admissible parent sets. Consider an arbitrary tuple
\[
(P_1,\ldots,P_p)
\in
\mathcal P_1\times\cdots\times\mathcal P_p,
\]
and let \(G(P_1,\ldots,P_p)\) be the directed graph obtained by assigning
\(P_j\) as the parent set of \(X_j\). If
\(X_k\rightarrow X_j\) belongs to this graph, admissibility implies
\(\rho(k)\prec\rho(j)\). Hence the role index strictly increases along every
directed path. A directed cycle would require nodes
\(X_{j_1},\ldots,X_{j_m}\) satisfying
\[
\rho(j_1)
\prec
\rho(j_2)
\prec
\cdots
\prec
\rho(j_m)
\prec
\rho(j_1),
\]
which is impossible under a strict ordering. Every tuple of admissible parent
sets therefore defines a DAG. Conversely, each graph in
\(\mathcal G_{\mathrm{role}}\) is uniquely determined by the tuple formed by
its node-specific parent sets. The admissible graph space is consequently in
one-to-one correspondence with
\(\mathcal P_1\times\cdots\times\mathcal P_p\).

Using the decomposability of BIC and the absence of any remaining graph-level
constraint coupling the parent-set choices,
\[
\begin{aligned}
\min_{G\in\mathcal G_{\mathrm{role}}}
\operatorname{BIC}(G)
&=
\min_{\substack{
(P_1,\ldots,P_p)\in
\mathcal P_1\times\cdots\times\mathcal P_p}}
\sum_{j=1}^{p}
\operatorname{BIC}_j(P_j)
\\
&=
\sum_{j=1}^{p}
\min_{P\in\mathcal P_j}
\operatorname{BIC}_j(P).
\end{aligned}
\]
Therefore, for any choices
\[
\widehat P_j
\in
\arg\min_{P\in\mathcal P_j}
\operatorname{BIC}_j(P),
\qquad
j=1,\ldots,p,
\]
the graph
\(\widehat G=G(\widehat P_1,\ldots,\widehat P_p)\)
minimizes the BIC over
\(\mathcal G_{\mathrm{role}}\).
\end{proof}

The result establishes global optimality only within the predefined
role-constrained graph class. It does not imply optimality over the
unrestricted space of DAGs. The argument also relies on the absence of
additional global restrictions coupling different nodes. For example, a
constraint on the total number of edges in the complete network would prevent
the local parent-set problems from being optimized independently.

Structure learning can therefore be implemented by exhaustive local
enumeration. For each node \(X_j\), every candidate
\(P\in\mathcal P_j\) is evaluated once and a BIC-minimizing parent set is
retained. The total number of local score evaluations is

\[
N_{\mathrm{scores}}
=
\sum_{j=1}^{p}
\sum_{d=0}^{\min(d_{\max},|\mathcal U_j|)}
\binom{|\mathcal U_j|}{d}.
\]

Thus, the substantive restrictions introduced from the accident-process roles
also have a computational consequence: they transform the constrained DAG
optimization problem into a collection of finite, independent parent-set
optimizations. Unlike greedy hill climbing, the resulting structure does not
depend on an initial graph, an edge-update sequence, or a local stopping rule.

If several parent sets achieve exactly the same minimum local BIC, ties are
resolved deterministically by retaining the parent set with the smallest
cardinality and, if required, by a fixed ordering of factor indices. This
rule affects only the choice among score-equivalent solutions and ensures
reproducibility of the learned structure.


% ------------------------------------------------------------
\subsubsection{Conditional-probability estimation and support diagnostics}
\label{sec:chap4_global_bn_parameters}
% ------------------------------------------------------------

Maximum-likelihood estimates are used in the BIC calculations above because
the criterion is defined from the maximized likelihood. After the structure
\(\widehat G\) has been selected, its conditional-probability tables are
first assessed on the complete accident sample. For every observed parent
configuration \(\pi\), the maximum-likelihood estimate is
\[
\widehat\theta_{j,\pi}
=
\frac{N_{j,\pi,1}}{N_{j,\pi}}.
\]
Parent configurations with \(N_{j,\pi}=0\) do not contribute to the local
log-likelihood or BIC during structure learning.

Before using the fitted network for joint-probability calculations, the
empirical support of every selected CPT row is recorded. The diagnostic
reports the numbers of rows with \(N_{j,\pi}=0\),
\(N_{j,\pi}\leq2\), and \(N_{j,\pi}\leq5\), the number of boundary MLEs
equal to zero or one, and the minimum and median observed cell counts. If all
rows required by the selected network meet the prespecified support criterion,
the final conditional-probability tables remain the MLEs. In that case the
fitted joint distribution is denoted by \(P_{\widehat G,\widehat\Theta}\).

The final estimator used for probabilistic inference is reported together
with these support diagnostics. If the diagnostic requires post-selection
regularization, its specification and an MLE-versus-regularized sensitivity
analysis of model-implied scenario probabilities are reported with the
results. This choice does not alter the BIC scores or the exact-optimization
result.


% ------------------------------------------------------------
\subsubsection{Bootstrap assessment of structural stability}
\label{sec:chap4_global_bn_bootstrap}
% ------------------------------------------------------------

Exact optimization removes uncertainty associated with the use of a local
search algorithm, but it does not remove sampling variability. The selected
structure may change when the set of observed accident narratives changes.
Structural stability is therefore assessed by nonparametric bootstrap
resampling at the accident level, following the use of bootstrap confidence
measures for features of learned Bayesian networks
\citep{friedman1999data}.

The semantic factors and their factual-unit assignments are kept fixed
throughout this analysis. The bootstrap therefore evaluates the stability of
the accident-level Bayesian-network structure conditional on the semantic
representation established in the preceding stage.

For each bootstrap replicate
\(r=1,\ldots,R_{\mathrm{boot}}\),
\(N_A\) accident narratives are sampled with replacement from the original
accident-level sample, and the exact structure-learning procedure is repeated
using the same role restrictions, maximum indegree and BIC criterion. Let
\(\widehat G^{[r]}\) denote the resulting graph. For an admissible dependency
\(X_a\rightarrow X_b\), its bootstrap selection frequency is

\[
\widehat f_{a\rightarrow b}
=
\frac{1}{R_{\mathrm{boot}}}
\sum_{r=1}^{R_{\mathrm{boot}}}
\mathbb I
\left\{
X_a\rightarrow X_b
\in
\widehat G^{[r]}
\right\}.
\]

The quantity
\(\widehat f_{a\rightarrow b}\)
measures how reproducibly the corresponding dependency is selected when the
accident sample is perturbed. Since edge orientation is fixed by the
accident-process roles, it concerns stability of edge selection under the
imposed direction rather than uncertainty between alternative orientations.
It is not interpreted as a posterior edge probability, an edge-wise
significance level, or evidence of a causal relationship.

The full-sample graph \(\widehat G\) remains the primary fitted structure.
Bootstrap frequencies complement its interpretation and may be used to
highlight dependencies exceeding a predefined descriptive stability threshold
\(\tau_{\mathrm{stab}}\) in reduced graphical representations. This threshold
does not modify the full-sample network.


% ------------------------------------------------------------
\subsubsection{Direction and magnitude of learned conditional dependencies}
\label{sec:chap4_global_bn_contrast}
% ------------------------------------------------------------

Selection of an edge \(X_a\rightarrow X_b\) indicates that the corresponding
parent is retained in a BIC-optimal local model of \(X_b\) within the
constrained graph class. Edge presence alone, however, does not indicate
whether \(X_a\) is associated with a higher or lower conditional probability
of \(X_b\), nor does it quantify the magnitude of that association.

For binary variables, the direction can be characterized through
context-specific conditional probability differences. This construction is
closely related to the notion of qualitative probabilistic influence
introduced by \citet{wellman1990fundamental}, in which the sign of a
probabilistic relationship is assessed while holding the remaining relevant
variables fixed.

For edge \(X_a\rightarrow X_b\), let

\[
\mathcal Q_{ab}
=
\operatorname{Pa}_{\widehat G}(b)
\setminus\{a\}
\]

denote the remaining direct parents of \(X_b\). For configuration \(\pi\) of
these variables, the fitted conditional probability contrast is

\[
\widetilde{\delta}_{a\rightarrow b}(\pi)
=
\widetilde P
\left(
X_b=1
\mid
X_a=1,
\mathbf X_{\mathcal Q_{ab}}=\pi
\right)
-
\widetilde P
\left(
X_b=1
\mid
X_a=0,
\mathbf X_{\mathcal Q_{ab}}=\pi
\right),
\]

where the probabilities are obtained from the post-selection conditional
tables defined above. For a fixed configuration of the other parents, a
positive contrast indicates a higher fitted probability of identifying
\(X_b\) when \(X_a\) is identified, whereas a negative contrast indicates the
opposite relationship.

For interpretation of the configuration-specific contrasts, attention is
restricted to contexts in which both states of \(X_a\) are represented in the
observed accident sample. Let

\[
\mathcal S_{ab}
=
\left\{
\pi:
N_{ab,1}(\pi)>0
\text{ and }
N_{ab,0}(\pi)>0
\right\}
\]

denote this empirically supported set of contexts. A dependency is described
as monotonically positive over the observed support when
\(\widetilde{\delta}_{a\rightarrow b}(\pi)\geq0\)
for every \(\pi\in\mathcal S_{ab}\), with at least one strict inequality.
It is described as monotonically negative under the reverse condition and as
non-monotone when contrasts of both signs occur.

To summarize the fitted association over the observed contexts, let

\[
\widehat w_{ab}(\pi)
=
\frac{
N_{ab}(\pi)
}{
\sum_{\pi'\in\mathcal S_{ab}}
N_{ab}(\pi')
}
\]

denote the empirical frequency of each estimable configuration of the
remaining parents. The observed-context weighted average contrast is

\[
\widetilde{\Delta}_{a\rightarrow b}
=
\sum_{\pi\in\mathcal S_{ab}}
\widehat w_{ab}(\pi)
\widetilde{\delta}_{a\rightarrow b}(\pi).
\]

The configuration-specific contrasts and their weighted average are retained
together. A positive
\(\widetilde{\Delta}_{a\rightarrow b}\)
indicates a positive conditional association on average over the empirically
supported contexts, but does not imply positivity in every context.

Bootstrap selection frequency and conditional probability contrast therefore
describe different properties. The former measures reproducibility of
structure selection, whereas the latter characterizes the fitted direction
and magnitude of a selected dependency. A dependency may be repeatedly
selected while having a small probability contrast, or display a larger
contrast while being less stable under resampling. Neither quantity is
interpreted causally. In prevention terms, these quantities help distinguish
factors that repeatedly belong to the same accident configuration from
relationships that remain directly informative after conditioning on the
other selected upstream factors.


% ============================================================
\subsection{Integration of recurrent scenarios and Bayesian dependencies}
\label{sec:chap4_scenario_bn_integration}
% ============================================================

Scenario mining and Bayesian-network modeling are integrated only after both
analyses have been estimated independently. This design preserves the
distinction between recurrence, association and conditional dependence.
Scenario support measures how frequently the complete accident configuration
occurs; confidence and lift characterize association between its antecedent
and consequence; and the Bayesian network characterizes the conditional
organization of the factor variables under the predefined role constraints.

The integration addresses two additional questions. First, which
relationships among the factors composing a recurrent scenario remain
represented as direct conditional dependencies in the full network? Second,
to what extent does the joint probability distribution implied by the network
reproduce the empirical frequency of the complete scenario?


% ------------------------------------------------------------
\subsubsection{Conditional structure within recurrent scenarios}
\label{sec:chap4_scenario_internal_bn}
% ------------------------------------------------------------

For retained scenario \(s\), define

\[
\widehat E_s
=
\left\{
X_a\rightarrow X_b\in\widehat E:
a\in\mathcal S_s,\;
b\in\mathcal S_s
\right\},
\]

where
\(\widehat G=(V,\widehat E)\)
is the network learned from the full accident sample. These internal edges
identify relationships among the factors of the scenario that remain directly
represented after conditioning on the other selected parents in the complete
network.

Particular attention is given to the two accident-process stages linking the
upstream A0/A1 factors to event \(b_s\) and event \(b_s\) to consequence
\(c_s\). For each internal dependency, its bootstrap selection frequency
\(\widehat f_{a\rightarrow b}\), its weighted conditional contrast
\(\widetilde{\Delta}_{a\rightarrow b}\), and, where relevant, the
configuration-specific contrasts are examined jointly. Dependencies exceeding
the descriptive stability threshold and exhibiting a positive conditional
contrast provide the strongest network representation of a positive
relationship between the corresponding scenario components.

The presence or absence of such edges does not validate or invalidate the
scenario itself. A recurrent scenario is defined by repeated joint occurrence
and may remain empirically meaningful even when some of its components are
not directly connected in the Bayesian network. Conversely, a stable
conditional dependency may occur without belonging to a complete scenario
that satisfies the recurrence criterion.

This distinction has a direct prevention-oriented interpretation. Scenario
mining identifies combinations of work context, adverse conditions, events
and consequences that repeatedly characterize accidents. The Bayesian network
then indicates which relationships within those configurations remain
informative once the other modeled parent factors are taken into account.
It therefore helps distinguish the overall recurrent context from the more
specific conditional organization represented by the fitted model.


% ------------------------------------------------------------
\subsubsection{Scenario recurrence relative to the Bayesian background model}
\label{sec:chap4_scenario_bn_expected}
% ------------------------------------------------------------

Analysis of internal network edges remains local to the relationships between
scenario components and does not determine whether the fitted network
reproduces the frequency of the complete configuration. Bayesian networks
have previously been used as background-knowledge models for frequent-pattern
analysis by comparing the empirical support of an itemset with the support
implied by the network
\citep{jaroszewicz2004interestingness,
jaroszewicz2009scalable}. We adopt this support-comparison principle after
recurrent scenarios have been identified, without using the Bayesian network
to determine which configurations satisfy the recurrence criterion.

For scenario \(\mathcal S_s\), the observed probability is its empirical
support,

\[
\widehat p_{\mathrm{obs},s}
=
\frac{n_s}{N_A}.
\]

The fitted Bayesian network defines a joint probability distribution over the
complete factor vector. The support implied by this distribution for the same
scenario is

\[
\widehat p_{\mathrm{BN},s}
=
P_{\widehat G,\widetilde\Theta}
\left(
X_g=1
\text{ for every }
g\in\mathcal S_s
\right).
\]

Since empirical scenario occurrence is defined by containment rather than
exact equality with the complete factor profile, all factors outside
\(\mathcal S_s\) must be marginalized. Equivalently,

\[
\widehat p_{\mathrm{BN},s}
=
\sum_{\substack{
\mathbf x\in\{0,1\}^{p}\\
x_g=1,\;g\in\mathcal S_s}}
\prod_{j=1}^{p}
P_{\widehat G,\widetilde\Theta}
\left(
X_j=x_j
\mid
\mathbf X_{\operatorname{Pa}_{\widehat G}(j)}
=
\mathbf x_{\operatorname{Pa}_{\widehat G}(j)}
\right).
\]

In implementation, this marginal probability is obtained using exact
Bayesian-network inference rather than explicit enumeration of all states of
the remaining variables.

The corresponding expected number of scenario occurrences under the fitted
network is

\[
\widehat\mu_s
=
N_A\widehat p_{\mathrm{BN},s}.
\]

Following \citet{jaroszewicz2004interestingness}, the model-relative
interestingness of the scenario can be expressed through the absolute
difference between empirical and BN-implied support,

\[
\widehat I_{\mathrm{BN},s}
=
\left|
\widehat p_{\mathrm{obs},s}
-
\widehat p_{\mathrm{BN},s}
\right|.
\]

Because the direction of this discrepancy is relevant for the present
application, the corresponding signed difference is additionally retained,

\[
\widehat D_{\mathrm{BN},s}
=
\widehat p_{\mathrm{obs},s}
-
\widehat p_{\mathrm{BN},s}.
\]

The absolute quantity
\(\widehat I_{\mathrm{BN},s}\)
measures the magnitude of disagreement between empirical scenario support and
the fitted Bayesian background model. The signed quantity distinguishes
configurations that occur more frequently than implied by the network
(\(\widehat D_{\mathrm{BN},s}>0\)) from those occurring less frequently than
the network assigns them
(\(\widehat D_{\mathrm{BN},s}<0\)).
On the accident-count scale, the same discrepancy is
\(n_s-\widehat\mu_s
=
N_A\widehat D_{\mathrm{BN},s}\).

These quantities are descriptive model-discrepancy measures. They are not
treated as p-values or as tests of excess recurrence, particularly because
the same accident sample contributes both to network estimation and to
empirical scenario support. Their purpose is instead to determine whether a
recurrent configuration is well represented by the conditional structure of
the fitted network.

For prevention, a recurrent scenario with empirical support close to its
BN-implied support is compatible with the probabilistic organization already
represented by the model. A scenario with a substantial positive discrepancy
is different: its complete configuration occurs more frequently than the
fitted network represents. Such a discrepancy does not establish an omitted
causal mechanism, but it identifies a configuration for which further expert
investigation may be informative. Relevant explanations may include an
unrepresented work practice, an organizational condition, an interaction not
captured by the permitted graph structure, a relevant omitted factor, or a
systematic feature of accident reporting.


% ------------------------------------------------------------
\subsubsection{Complementary interpretation for prevention}
\label{sec:chap4_scenario_interpretation}
% ------------------------------------------------------------

The final analysis deliberately preserves the distinction between four
properties of a recurrent accident scenario. Its support describes empirical
recurrence. Confidence and lift describe association between its
upstream/event antecedent and its consequence. Internal Bayesian-network
edges and conditional probability contrasts characterize direct conditional
dependencies under the imposed accident-process organization, while bootstrap
frequencies assess the reproducibility of those selected dependencies.
Finally, comparison between empirical and BN-implied support characterizes
the extent to which the fitted joint model reproduces the recurrence of the
complete configuration.

These quantities are not combined into a single score because they answer
different questions. A highly recurrent scenario need not be strongly
enriched for a particular consequence; a large conditional association need
not correspond to a recurrent complete configuration; and two scenarios with
similar empirical recurrence may differ substantially in their discrepancy
relative to the Bayesian-network background model.

This distinction is useful for prevention because it separates three levels
of interpretation: the configurations that repeatedly occur, the
relationships within those configurations that remain informative after
conditioning, and the recurrent configurations that are not well represented
by the current probabilistic model. The last group provides targets for
additional expert investigation rather than automatic prevention priorities.
Throughout the analysis, recurrence, association, conditional dependence,
structural stability and model discrepancy remain distinct descriptive
properties and none is interpreted as evidence of a causal effect.
